% model-ru.tex — математическая модель движения первой ступени РН
% на участках выведения, возвращения и посадки.
% Сборка: docs\build-ru.ps1   (или pdflatex model-ru.tex дважды)
\documentclass[11pt,a4paper]{article}

\usepackage[T2A]{fontenc}
\usepackage[utf8]{inputenc}
\usepackage{lmodern}
\usepackage[russian]{babel}
\usepackage{amsmath,amssymb,bm}
\usepackage{booktabs}
\usepackage{geometry}
\geometry{margin=2.2cm}
\usepackage[hidelinks]{hyperref}

\newcommand{\vect}[1]{\bm{#1}}
\newcommand{\quat}[1]{\bm{#1}}
\newcommand{\Rot}{\mathcal{R}}
\newcommand{\dd}{\mathrm{d}}
\newcommand{\eh}{\hat{\vect{e}}_d}
\newcommand{\ec}{\hat{\vect{e}}_c}

\title{Математическая модель движения первой ступени\\
ракеты-носителя на участках выведения,\\
возвращения и посадки}
\author{}
\date{}

\begin{document}
\maketitle
\vspace{-2.4cm}

\begin{abstract}
\noindent
Изложена математическая модель движения первой ступени ракеты-носителя,
совершающей автономный полёт по схеме возвращения к месту старта. Модель
охватывает участок выведения, выключение маршевых двигателей, импульс
возврата, баллистический участок с прохождением вершины траектории, тормозной
импульс входа в атмосферу, аэродинамический спуск на решётчатых рулях,
посадочный тормозной импульс и контактное взаимодействие опор с посадочной
поверхностью. Приведены уравнения движения, модель стандартной атмосферы,
аэродинамические характеристики ступени в двух положениях корпуса, программа
управления на участке выведения, постановка задачи возврата как краевой,
алгоритмы наведения атмосферного участка, модель контакта, а также
алгоритм наведения посадочного участка, построенный методом последовательной
выпуклой аппроксимации. Указаны принятые допущения.

\medskip\noindent
\textbf{Ключевые слова:} ракета-носитель, многоразовая первая ступень,
возвращение к месту старта, импульс возврата, реактивное торможение,
решётчатые рули, посадочное устройство, математическое моделирование
\end{abstract}

\tableofcontents

% =====================================================================
\section{Введение}

Многоразовая первая ступень ракеты-носителя (РН), возвращаемая к месту старта,
выполняет за один полёт несколько качественно различных манёвров. Участок
выведения проходит в плотных слоях атмосферы при работающей двигательной
установке; импульс возврата и последующий баллистический участок — практически
в вакууме; вход в атмосферу, аэродинамический спуск и посадка — снова в плотных
слоях, но при обратной ориентации корпуса. Разнородность этих режимов не
позволяет описать весь полёт единой системой уравнений приемлемой сложности.

Принято следующее разделение. Внеатмосферный участок — от старта до пересечения
\emph{границы входа} $h_{\text{гр}}=65$~км на снижении — описан плоской моделью
точечной массы над сферической невращающейся Землёй с учётом аэродинамического
сопротивления на участке выведения. Атмосферный участок, расположенный ниже
границы входа, описан пространственной моделью твёрдого тела с шестью степенями
свободы, аэродинамикой корпуса и решётчатых рулей и моделью контакта посадочных
опор. Сопряжение участков выполняется по вектору состояния; уравнения движения
двух моделей не смешиваются.

Правомерность такого разделения обоснована малостью аэродинамических сил выше
границы входа. При скорости $1579$~м/с скоростной напор составляет $386$~Па на
высоте 60~км, $23$~Па на высоте 80~км и $0{,}70$~Па на высоте 100~км, что для
ступени массой 80~т даёт тормозное ускорение порядка
$8\cdot10^{-5}\ \text{м/с}^2$, то есть около четырёх стотысячных долей от
ускорения, создаваемого тягой.

Далее приведены все уравнения и алгоритмы в том виде, в каком они реализованы в
расчётной программе.

% =====================================================================
\section{Расчётная схема и системы координат}

Инерциальная система координат — правая, ось $\hat{\vect{z}}$ направлена вверх,
начало совмещено с точкой старта, она же точка посадки. Связанная система
координат имеет ось $\hat{\vect{z}}_B$, направленную вдоль продольной оси
ступени от двигательной установки к носовой части.

Ориентация задаётся кватернионом со скалярной частью впереди,
$\quat{q}=(q_0,q_1,q_2,q_3)$, $\lVert\quat{q}\rVert=1$, переводящим связанную
систему координат в инерциальную:
\begin{equation}
\vect{x}_I=\Rot(\quat{q})\,\vect{x}_B .
\end{equation}
Символ $\otimes$ обозначает произведение кватернионов. Ускорение свободного
падения на поверхности $g_0=9{,}80665\ \text{м/с}^2$. Всюду применяется система
единиц СИ.

Угол отклонения продольной оси от местной вертикали определяется кватернионом:
\begin{equation}
\cos\theta = 1-2\left(q_1^2+q_2^2\right).
\label{eq:tilt}
\end{equation}

% =====================================================================
\section{Модель атмосферы}

Принята стандартная атмосфера, заданная на двух участках по высоте.

Ниже геометрической высоты 86~км применяется семислойная модель с кусочно-линейным
профилем температуры, вычисляемая по геопотенциальной высоте
\begin{equation}
H = \frac{R' h}{R'+h},
\qquad R'=6356{,}766\ \text{км},
\end{equation}
\begin{equation}
T = T_b + L_b (H - H_b),
\qquad
p = \begin{cases}
p_b\left(\dfrac{T_b}{T}\right)^{g_0/(R_{\text{в}}L_b)}, & L_b \neq 0,\\[2ex]
p_b\exp\!\left(-\dfrac{g_0 (H-H_b)}{R_{\text{в}}T_b}\right), & L_b = 0,
\end{cases}
\qquad \rho = \frac{p}{R_{\text{в}}T},
\end{equation}
где $R_{\text{в}}=287{,}053\ \text{Дж}/(\text{кг}\cdot\text{К})$ — удельная
газовая постоянная воздуха. Скорость звука $a=\sqrt{\gamma R_{\text{в}}T}$,
показатель адиабаты $\gamma=1{,}4$. Параметры базовых слоёв приведены в
табл.~\ref{tab:atm}.

\begin{table}[htb]
\centering
\caption{Параметры базовых слоёв модели атмосферы}
\label{tab:atm}
\begin{tabular}{lrrrrrrr}
\toprule
$H_b$, км & 0 & 11 & 20 & 32 & 47 & 51 & 71 \\
\midrule
$L_b$, К/м & $-0{,}0065$ & 0 & $0{,}001$ & $0{,}0028$ & 0 & $-0{,}0028$ & $-0{,}002$ \\
$T_b$, К & 288{,}15 & 216{,}65 & 216{,}65 & 228{,}65 & 270{,}65 & 270{,}65 & 214{,}65 \\
$p_b$, Па & 101325 & 22632{,}06 & 5474{,}889 & 868{,}019 & 110{,}906 & 66{,}939 & 3{,}956 \\
\bottomrule
\end{tabular}
\end{table}

Выше 86~км кусочно-линейное представление профиля температуры неприменимо
вследствие изменения молекулярной массы воздуха, поэтому плотность определяется
логарифмической интерполяцией табличных значений (табл.~\ref{tab:atm2}).
Использование вместо этого экспоненциальной экстраполяции завышает плотность на
высоте 100~км примерно на 20\,\%, что существенно, поскольку именно малость
плотности на этих высотах обосновывает принятое разделение участков полёта.
Отметим, что геопотенциальная высота $84{,}852$~км в точности соответствует
геометрической высоте 86~км, вследствие чего оба участка модели стыкуются без
разрыва.

\begin{table}[htb]
\centering
\caption{Плотность воздуха на высотах свыше 86 км}
\label{tab:atm2}
\begin{tabular}{lrrrrrrr}
\toprule
$h$, км & 86 & 90 & 95 & 100 & 110 & 120 & 130 \\
\midrule
$\rho$, кг/м$^3$ & $6{,}958{\cdot}10^{-6}$ & $3{,}416{\cdot}10^{-6}$ &
$1{,}393{\cdot}10^{-6}$ & $5{,}604{\cdot}10^{-7}$ & $9{,}708{\cdot}10^{-8}$ &
$2{,}222{\cdot}10^{-8}$ & $8{,}152{\cdot}10^{-9}$ \\
\bottomrule
\end{tabular}
\end{table}

Реализованная модель воспроизводит табличные значения плотности с погрешностью
не более $0{,}01\,\%$ в диапазоне высот от уровня моря до 120~км.

% =====================================================================
\section{Аэродинамические характеристики ступени}
\label{sec:aero}

Ступень обтекается в двух существенно различных положениях, которым
соответствуют различные аэродинамические коэффициенты.

На участке выведения ступень движется носовой частью вперёд. Открытый отсек
межступенного перехода закрыт носовым обтекателем — конусом длиной $6{,}0$~м с
полууглом $17^\circ$ при диаметре корпуса $3{,}66$~м. Для заострённого
конус-цилиндра при нулевом угле атаки коэффициент лобового сопротивления
составляет около $0{,}30$ на дозвуковых скоростях, достигает $0{,}55$ в
трансзвуковой области и снижается до $0{,}24$ при $M=8$.

На участке возвращения ступень движется двигательной установкой вперёд. Здесь
передней является плоское донное сечение с развитым отрывом потока, то есть
плохообтекаемое тело, и коэффициент осевой силы существенно выше.

Обе зависимости заданы таблично (табл.~\ref{tab:cd}) и отнесены к площади миделя
$S=\pi(D/2)^2=10{,}52\ \text{м}^2$. Произведение $C_D S$ составляет
соответственно $3{,}2\ldots5{,}79\ldots2{,}5\ \text{м}^2$ и
$8{,}94\ldots14{,}73\ \text{м}^2$, то есть различается примерно в два с половиной
раза. Это различие определяется положением корпуса относительно набегающего
потока, а не неопределённостью исходных данных, и его необходимо учитывать при
сопоставлении с литературными значениями.

\begin{table}[htb]
\centering
\caption{Аэродинамические коэффициенты в зависимости от числа Маха}
\label{tab:cd}
\begin{tabular}{lrrrrrrrrrrr}
\toprule
\multicolumn{12}{l}{\emph{Носовой частью вперёд (участок выведения)}}\\
$M$ & 0 & 0{,}6 & 0{,}8 & 0{,}9 & 1{,}05 & 1{,}2 & 1{,}5 & 2{,}0 & 3{,}0 & 5{,}0 & 8{,}0 \\
$C_D$ & 0{,}30 & 0{,}30 & 0{,}32 & 0{,}38 & 0{,}55 & 0{,}52 & 0{,}45 & 0{,}38 & 0{,}30 & 0{,}26 & 0{,}24 \\
\midrule
\multicolumn{12}{l}{\emph{Двигательной установкой вперёд (участок возвращения)}}\\
$M$ & 0 & 0{,}6 & 0{,}9 & 1{,}1 & 1{,}5 & 2{,}0 & 3{,}0 & 5{,}0 & 8{,}0 & & \\
$C_A$ & 0{,}85 & 0{,}88 & 1{,}10 & 1{,}40 & 1{,}25 & 1{,}10 & 0{,}95 & 0{,}92 & 0{,}90 & & \\
\midrule
\multicolumn{12}{l}{\emph{Производная нормальной силы по углу атаки, 1/рад}}\\
$C_{N\alpha}$ & 2{,}0 & 2{,}1 & 2{,}4 & 2{,}5 & 2{,}4 & 2{,}3 & 2{,}2 & 2{,}1 & 2{,}0 & & \\
\bottomrule
\end{tabular}
\end{table}

% =====================================================================
\section{Внеатмосферный участок: уравнения движения}

Движение рассматривается в плоскости, в геоцентрической невращающейся системе
координат; точка старта имеет координаты $(0,R_\oplus)$, $R_\oplus=6371$~км.
Вектор состояния $\vect{y}=(X,Z,V_X,V_Z,m)$. Поле тяготения центральное,
гравитационный параметр выбран из условия точного равенства ускорения
свободного падения на поверхности величине $g_0$:
\begin{equation}
\mu = g_0 R_\oplus^2,
\qquad
\vect{a}_g = -\frac{\mu}{\lVert\vect{r}\rVert^3}\,\vect{r} .
\end{equation}
Высота, дальность по трассе и орты местной вертикали и направления полёта:
\begin{equation}
h = \lVert\vect{r}\rVert - R_\oplus,
\qquad
d = R_\oplus \arctan\frac{X}{Z},
\qquad
\hat{\vect{u}} = \frac{\vect{r}}{\lVert\vect{r}\rVert},
\qquad
\hat{\vect{e}} = \frac{(Z,-X)}{\lVert\vect{r}\rVert} .
\end{equation}
Учёт кривизны Земли существен: на дальности 70~км поверхность понижается на
384~м.

Уравнения движения:
\begin{equation}
\dot{\vect{v}} = \vect{a}_g
  + \frac{n_e\,\tau\,T_e(p)}{m}\,\hat{\vect{u}}_T
  - \frac{q S C_D(M)}{m}\,\frac{\vect{v}}{\lVert\vect{v}\rVert},
\qquad
\dot m = -n_e\,\tau\,\dot m_e ,
\end{equation}
где $q=\tfrac12\rho\lVert\vect{v}\rVert^2$ — скоростной напор,
$M=\lVert\vect{v}\rVert/a$ — число Маха, $n_e$ — число работающих камер, $\tau$
— степень дросселирования, $\hat{\vect{u}}_T$ — орт направления тяги. Тяга
одной камеры с учётом противодавления:
\begin{equation}
T_e(p)=T_{\text{п}}-A_e p,
\qquad
\dot m_e = \frac{T_{\text{з}}}{g_0 I_{\text{з}}},
\qquad
T_{\text{п}} = \dot m_e\, g_0 I_{\text{п}},
\qquad
A_e = \frac{T_{\text{п}}-T_{\text{з}}}{p_0},
\end{equation}
где индексы «з» и «п» относятся соответственно к земным и пустотным условиям.

Интегрирование выполняется методом Рунге—Кутты четвёртого порядка с постоянным
шагом. Условие останова (достижение заданной скорости выключения, пересечение
заданной высоты, истечение времени работы двигателей) уточняется делением
последнего шага пополам, благодаря чему момент останова совпадает с условием, а
не с ближайшим последующим узлом сетки.

При $m\le m_{\text{сух}}$ тяга и секундный расход полагаются равными нулю. Без
этого условия вычислительный алгоритм, превысив располагаемый запас топлива,
продолжает интегрирование при отрицательной массе топлива, где отношение $T/m$
претерпевает разрыв.

% =====================================================================
\section{Участок выведения}

\subsection{Программа управления}

Программа угла тангажа состоит из вертикального подъёма, плавного завала по
сглаженной ступенчатой функции и последующего гравитационного разворота при
нулевом угле атаки. Постоянное смещение $\beta$ направления тяги от вектора
скорости допускается только выше высоты $h_{\text{св}}=40$~км, где скоростной
напор не превышает 300~Па и управление не приводит к росту поперечных нагрузок:
\begin{equation}
\hat{\vect{u}}_T(t) =
\begin{cases}
\hat{\vect{u}}, & t<t_p,\\[1ex]
\cos(\psi\tilde s)\,\hat{\vect{u}} + \sin(\psi\tilde s)\,\hat{\vect{e}},
  \quad \tilde s = s^2(3-2s),\ \ s=\dfrac{t-t_p}{t_r}, & t_p \le t < t_p+t_r,\\[2ex]
\hat{\vect{v}}, & h<h_{\text{св}},\\[1ex]
\cos\beta\,\hat{\vect{v}} + \sin\beta\,\hat{\vect{n}}, & h\ge h_{\text{св}},
\end{cases}
\end{equation}
где $t_p$ — момент начала завала, $t_r$ — его длительность, $\psi$ — угол
завала, $\hat{\vect{n}}\perp\hat{\vect{v}}$ — орт, лежащий в плоскости движения
со стороны возрастания дальности.

Дросселирование обеспечивает ограничение осевой перегрузки:
\begin{equation}
\tau = \min\!\left(1,\ \max\!\left(\tau_{\min},\
 \frac{n_{\text{пр}} g_0 m}{n_e T_e(p)}\right)\right).
\label{eq:throttle}
\end{equation}

Число работающих камер выбирается как наименьшее, обеспечивающее стартовую
тяговооружённость не ниже заданной:
\begin{equation}
n_e=\min\left\{n:\ n\,T_{\text{з}}\ge k_{\text{тв}}\,m_0 g_0\right\},
\qquad k_{\text{тв}}=1{,}35 .
\end{equation}

\subsection{Постановка задачи выведения}

Ступень, совершающая автономный полёт, не несёт полезной нагрузки и верхних
ступеней, вследствие чего её тяговооружённость существенно выше, чем у ступени в
составе носителя. При работе двигателей на номинальном режиме заданная скорость
выключения достигается на высоте около 54~км, то есть на $6{,}5$~км ниже
требуемой, причём смещение $\beta$ при этом достигает предельного значения.
Поэтому предельная осевая перегрузка $n_{\text{пр}}$ в \eqref{eq:throttle}
рассматривается не как ограничение, а как искомый параметр программы управления.

Определяемые параметры: начальная масса $m_0$, момент начала завала $t_p$, угол
завала $\psi$, смещение $\beta$ и предельная перегрузка $n_{\text{пр}}$. Условия,
проверяемые в момент достижения заданной скорости выключения $v_{\text{вык}}$:
высота $h_{\text{вык}}$, дальность $d_{\text{вык}}$ и масса $m_{\text{вык}}$.
Задача решается методом наименьших квадратов с доверительной областью;
оставшаяся степень свободы устраняется либо малым штрафом на $\beta$, что
отвечает программе с наименьшим управляющим воздействием, либо заданием угла
наклона траектории в момент выключения.

\subsection{Влияние угла наклона траектории в момент выключения}

Три условия выключения — высота, скорость и дальность — не определяют
направление вектора скорости. Между тем высота вершины траектории практически
полностью определяется его вертикальной составляющей. Результаты
параметрического исследования приведены в табл.~\ref{tab:gamma}.

\begin{table}[htb]
\centering
\caption{Зависимость характеристик полёта от угла наклона траектории
в момент выключения}
\label{tab:gamma}
\begin{tabular}{lrrrr}
\toprule
$\gamma_{\text{вык}}$, град & $v_x$, м/с & $m_{\text{т}}$, т
 & $h_{\max}$, км & $d_{\max}$, км \\
\midrule
$49{,}8$ (наименьшее управление) & 1020 & $151{,}4$ & $168{,}0$ & $60{,}8$ \\
$43{,}1$ & 1153 & $169{,}9$ & $154{,}7$ & $67{,}5$ \\
$42{,}4$ & 1167 & $170{,}8$ & $153{,}3$ & $68{,}3$ \\
\bottomrule
\end{tabular}

\smallskip
\footnotesize Примечание: $m_{\text{т}}$ — суммарный расход топлива на
выведение и манёвр возврата.
\end{table}

При $\gamma_{\text{вык}}$ ниже примерно $42^\circ$ смещение $\beta$ достигает
предельного значения и траектория становится нереализуемой при гравитационном
развороте с нулевым углом атаки. Программа с наименьшим управляющим
воздействием одновременно является наиболее экономичной и даёт наименьшую
максимальную дальность, вследствие чего принята в качестве основной.

% =====================================================================
\section{Манёвр возврата}

\subsection{Разворот ступени}

Вне атмосферы решётчатые рули не создают управляющего момента, а при выключенных
двигателях его не создаёт и отклонение камер. Управляющий момент реактивной
системы управления (РСУ) создаётся парой сопел:
\begin{equation}
M_{\text{РСУ}} = n_r F_r \left\lvert z_r - z_{\text{цм}}\right\rvert,
\qquad
\varepsilon_{\text{РСУ}} = \frac{M_{\text{РСУ}}}{J}
 \approx 0{,}31\ \text{град/с}^2 .
\end{equation}
Разворот на $150^\circ$ по релейному закону занял бы при этом $44$~с, в течение
которых ступень удаляется ещё на 49~км по дальности, что увеличивает
максимальную дальность до 97~км.

Поэтому принята следующая последовательность: осаждение топлива в течение
$t_{\text{ос}}=8$~с, запуск двигателей и завершение разворота отклонением камер.
Располагаемое угловое ускорение при работающих двигателях
\begin{equation}
\varepsilon_{\text{дв}}
 = \frac{n_r F_r\left\lvert z_r - z_{\text{цм}}\right\rvert
   + T\sin\delta_{\max}\,z_{\text{цм}}}{J}
 \approx 15{,}7\ \text{град/с}^2 ,
\end{equation}
то есть в пятьдесят раз больше, и тот же разворот занимает $5{,}7$~с.

Разворот и рабочий импульс совмещены: направление тяги изменяется по релейному
закону в течение горения, что учитывается при интегрировании. Угловой профиль:
\begin{equation}
t_{\text{р}} = 2\sqrt{\frac{\Delta}{\varepsilon_{\text{дв}}}},
\qquad
\theta(t)=\begin{cases}
\tfrac12\varepsilon_{\text{дв}} t^2, & t\le t_{\text{р}}/2,\\[1ex]
\Delta-\tfrac12\varepsilon_{\text{дв}}\left(t_{\text{р}}-t\right)^2,
 & t_{\text{р}}/2<t<t_{\text{р}},\\[1ex]
\Delta, & t\ge t_{\text{р}} .
\end{cases}
\end{equation}
Учёт совмещения разворота с рабочим импульсом уменьшает расчётную максимальную
дальность с 97 до 61~км.

\subsection{Краевая задача возврата}

Задача возврата является задачей приведения в заданную точку, а не задачей
мягкой посадки: в момент выключения импульса возврата ступень находится на
высоте около 100~км и движется со скоростью порядка километра в секунду.
Постановка следующая: по вектору состояния в момент выключения маршевых
двигателей найти ориентацию $\chi$, неизменную в инерциальной системе координат,
и длительность импульса $t_b$ такие, чтобы баллистическое продолжение траектории
приводило в точку старта.

Одно условие при двух неизвестных даёт однопараметрическое семейство решений;
оставшаяся степень свободы используется для минимизации расхода топлива, что при
неизменных числе камер и степени дросселирования эквивалентно минимизации
длительности импульса:
\begin{equation}
\min_{\chi\ \ge\ 180^\circ}\ t_b(\chi)
\qquad\text{при условии}\qquad
d_{\text{пад}}\left(\chi, t_b\right) = d_{\text{пр}} ,
\label{eq:bvp}
\end{equation}
где $d_{\text{пад}}$ — дальность точки падения баллистического продолжения,
$d_{\text{пр}}$ — дальность точки прицеливания.

Для каждого значения $\chi$ длительность определяется методом деления отрезка
пополам: невязка монотонна по $t_b$, поскольку увеличение длительности импульса
всегда смещает точку падения в сторону старта. Минимизация по $\chi$ выполняется
одномерным методом на отрезке. Верхняя граница отрезка по $t_b$ определяется
располагаемым запасом топлива:
\begin{equation}
t_b^{\max}=0{,}98\,\frac{m-m_{\text{сух}}}{n_e\dot m_e}.
\end{equation}

\subsection{Обоснование ограничения на ориентацию}

Ограничение $\chi\ge180^\circ$ в \eqref{eq:bvp} означает, что вектор тяги не
имеет составляющей, направленной вверх. Оно необходимо ввиду вырожденности
целевой функции: при изменении $\chi$ в пределах рабочего диапазона расход
топлива изменяется на $2\,\%$, тогда как высота вершины траектории и вертикальная
скорость на границе входа изменяются существенно (табл.~\ref{tab:chi}).

\begin{table}[htb]
\centering
\caption{Характеристики манёвра возврата при различной ориентации}
\label{tab:chi}
\begin{tabular}{rrrrr}
\toprule
$\chi$, град & $t_b$, с & $m_{\text{т}}$, т & $h_{\max}$, км
 & $v_z$ на границе входа, м/с \\
\midrule
$172{,}5$ & $35{,}60$ & $32{,}64$ & $163{,}4$ & $-1329$ \\
$180{,}0$ & $35{,}85$ & $32{,}86$ & $144{,}0$ & $-1184$ \\
$185{,}0$ & $36{,}32$ & $33{,}29$ & $132{,}5$ & $-1090$ \\
$195{,}0$ & $38{,}15$ & $34{,}97$ & $113{,}1$ & $-905$ \\
\bottomrule
\end{tabular}
\end{table}

Без ограничения решение смещается в область $\chi<180^\circ$, где ступень
получает дополнительную вертикальную скорость: высота вершины траектории
возрастает до 200~км, а скорость на границе входа — на 500~м/с при экономии
топлива около 30~кг. Целевая функция \eqref{eq:bvp} этого не учитывает,
поскольку определена только до момента выключения импульса возврата, тогда как
соответствующие затраты возникают позже — в тормозном импульсе входа.
Ограничение отражает физический смысл манёвра: импульс возврата является
тормозным, и сообщение ступени дополнительной энергии в вертикальном направлении
противоречит его назначению. Стоимость ограничения составляет $0{,}1\,\%$
расхода топлива.

% =====================================================================
\section{Баллистический участок и сопряжение}

\subsection{Баллистический участок}

Движение на баллистическом участке описывается теми же уравнениями при
выключенных двигателях; аэродинамические коэффициенты принимаются по
табл.~\ref{tab:cd} для положения двигательной установкой вперёд. Ориентация
приводится к противоположной вектору скорости силами РСУ; при
$\varepsilon_{\text{РСУ}}=0{,}31\ \text{град/с}^2$ и длительности участка около
пяти минут располагаемого управляющего момента достаточно с большим запасом.

\subsection{Коррекция точки прицеливания}

Прицеливание в \eqref{eq:bvp} выполняется по точке падения баллистического
продолжения в отсутствие атмосферы, что справедливо выше границы входа.
Действительное снижение происходит в атмосфере и укорачивает траекторию,
вследствие чего точка прицеливания корректируется по результату расчёта
атмосферного участка:
\begin{equation}
d_{\text{пр}}^{(k+1)} = d_{\text{пр}}^{(k)} - 0{,}9\,\epsilon^{(k)},
\end{equation}
где $\epsilon^{(k)}$ — отклонение точки касания от точки старта на $k$-й
итерации. Зависимость близка к линейной с единичным коэффициентом, вследствие
чего достаточно двух-трёх итераций.

\subsection{Замыкание баланса топлива}

Масса в момент выключения маршевых двигателей должна обеспечивать выполнение
манёвра возврата и последующих тормозных импульсов, тогда как расход на манёвр
возврата сам зависит от этой массы. Управляемой величиной является запас топлива
на границе входа $m_{\text{т,гр}}$. Производная
\begin{equation}
\frac{\partial m_{\text{т,гр}}}{\partial m_{\text{вык}}},
\end{equation}
вычисленная по формуле Циолковского при неизменной траектории, равна $0{,}63$,
тогда как её измеренное значение составляет $1{,}24$, поскольку при изменении
$m_{\text{вык}}$ заново определяется вся программа выведения. Ввиду этого
итерации выполняются методом секущих, а угол $\gamma_{\text{вык}}$ фиксируется
после первой итерации; в противном случае каждая итерация приводит к решению из
другой области притяжения, и процесс переходит в незатухающий двухтактный цикл.

\subsection{Выбор высоты границы входа}
\label{subsec:gr}

Тормозной импульс входа включается и выключается по прогнозируемому
максимальному скоростному напору (\S\ref{subsec:entry}). На высоте 70~км, где
$q=62$~Па, прогнозируемое значение опускается ниже порога выключения при остатке
топлива около 20~т, и импульс прекращается преждевременно. Оставшаяся масса
ступени составляет при этом 55~т, а посадочный тормозной импульс выполняется
одной камерой, что даёт результирующее ускорение и потребную высоту включения
\begin{equation}
a_{\text{р}} = \frac{T_{\text{з}}}{m} - g_0
 = \frac{845}{55{,}6} - 9{,}81 = 5{,}4\ \text{м/с}^2,
\qquad
h_{\text{вкл}} = \frac{v^2}{2a_{\text{р}}} = \frac{380^2}{10{,}8}
 = 13{,}4\ \text{км}.
\end{equation}
При таком включении запас топлива вырабатывается до момента касания.

При высоте границы входа ниже 66~км тормозной импульс входа продолжается до
достижения нижнего предела запаса топлива (7~т), масса ступени составляет 33~т, и
посадочный импульс включается на высоте $2{,}7$~км. Переход между режимами
происходит в диапазоне высот 66\dots68~км, вследствие чего граница входа принята
равной 65~км. Скоростной напор на этой высоте составляет около 130~Па, что даёт
тормозное ускорение $0{,}03\ \text{м/с}^2$ — величину на два порядка меньшую
границы применимости внеатмосферной постановки.

% =====================================================================
\section{Атмосферный участок: уравнения движения}

\subsection{Общий вид}

Движение описывается как движение твёрдого тела с шестью степенями свободы.
Вектор состояния включает положение $\vect{r}$, скорость $\vect{v}$, кватернион
ориентации $\quat{q}$, угловую скорость $\boldsymbol{\omega}$, массы компонентов
топлива, текущее значение дросселирования и углы отклонения рулей:
\begin{equation}
m\dot{\vect{v}} = \Rot(\quat{q})\left(\vect{F}_{\text{т}}
 +\vect{F}_{\text{а}}+\vect{F}_{\text{р}}\right)
 + m\vect{g} + \vect{F}_{\text{к}},
\qquad
\vect{J}(m)\dot{\boldsymbol{\omega}} = \vect{M}_B
  - \boldsymbol{\omega}\times\vect{J}(m)\boldsymbol{\omega},
\end{equation}
где $\vect{F}_{\text{т}}$, $\vect{F}_{\text{а}}$, $\vect{F}_{\text{р}}$,
$\vect{F}_{\text{к}}$ — силы тяги, аэродинамическая корпуса, аэродинамическая
рулей и контактная соответственно. Тензор инерции пересчитывается по текущим
массам компонентов. Шаг интегрирования $\Delta t=0{,}002$~с, частота работы
алгоритмов наведения 50~Гц.

\subsection{Массово-инерционные характеристики}

Сухая конструкция представлена тонкостенной оболочкой, топливо — сплошным
столбом; обе составляющие приводятся к общему центру масс по теореме
Гюйгенса—Штейнера. При массе топлива $m_{\text{т}}=\max(m-m_{\text{сух}},0)$,
площади поперечного сечения столба $A$, плотности $\rho_{\text{т}}$ и
$r_c^2=A/\pi$:
\begin{gather}
h_c = \frac{m_{\text{т}}}{\rho_{\text{т}} A},\qquad
z_c = z_{\text{осн}} + \tfrac{1}{2}h_c,\qquad
z_{\text{цм}} = \frac{m_{\text{сух}}z_{\text{сух}}
 + m_{\text{т}} z_c}{m},\\[4pt]
I_{\text{сух}} = \frac{m_{\text{сух}}\left(6R^2+L^2\right)}{12},
\qquad
I_{c} = \frac{m_{\text{т}}\left(3r_c^2+h_c^2\right)}{12},\\[4pt]
J = I_{\text{сух}}
 + m_{\text{сух}}\left(z_{\text{сух}}-z_{\text{цм}}\right)^2
 + I_{c} + m_{\text{т}}\left(z_c-z_{\text{цм}}\right)^2 .
\label{eq:inertia}
\end{gather}

\subsection{Тяга}

\begin{equation}
T = n_e\,\tau_{\text{д}}\left(T_{\text{п}}-A_e p\right)
 k_{\text{з}}(t_{\text{з}}),
\qquad
\vect{M}_{\text{т}} = \left(\vect{r}_{\text{к}}-\vect{r}_{\text{цм}}\right)
 \times\vect{F}_{\text{т}},
\end{equation}
где $\tau_{\text{д}}$ — действительное значение дросселирования, отличающееся от
командного вследствие ограничения скорости изменения ($2$~с$^{-1}$) и
апериодического запаздывания с постоянной времени $0{,}1$~с; $k_{\text{з}}$ —
характеристика выхода на режим, линейно возрастающая от нуля при
$t_{\text{з}}=0{,}5$~с до единицы при $t_{\text{з}}=1{,}5$~с; при выключении тяга
спадает по экспоненте с постоянной времени $0{,}15$~с; $\vect{r}_{\text{к}}$ —
точка приложения тяги.

\subsection{Аэродинамические силы и моменты корпуса}

При скорости относительно воздуха $\vect{v}_{\text{о}}=\vect{v}-\vect{w}$,
скоростном напоре $q=\tfrac12\rho v_{\text{о}}^2$ и угле атаки $\alpha$,
определяемом соотношением $\cos\alpha=-\hat v_{\text{о},z}$, осевая и нормальная
силы равны
\begin{equation}
F_{\text{ос}} = -\operatorname{sign}\left(\hat v_{\text{о},z}\right)
 q S C_A(M)\,k_{\text{э}},
\qquad
F_N = q S\, C_{N\alpha}(M)\,\alpha\, k_{\text{э}} .
\end{equation}
Нормальная сила приложена в центре давления
\begin{equation}
x_{\text{цд}} = f(M,\alpha)\,L_{\text{ст}},
\qquad
f(M,\alpha) = 0{,}28
 + 0{,}01\exp\!\left[-\left(\frac{M-1{,}05}{0{,}3}\right)^2\right]
 - 0{,}015\min\!\left(\frac{\lvert\alpha\rvert}{15^\circ},\,1\right).
\end{equation}
Поскольку $f<z_{\text{цм}}/L_{\text{ст}}\approx0{,}298$ во всём диапазоне чисел
Маха, корпус без рулей статически неустойчив; устойчивость обеспечивается
решётчатыми рулями.

Экранирование корпуса струёй при работающих двигателях учитывается
коэффициентом $k_{\text{э}}$: при коэффициенте тяги $C_T=T/(qS)>0{,}5$
\begin{equation}
k_{\text{э}} = 1 - 0{,}95\min\!\left(\frac{C_T-0{,}5}{2{,}5},\,1\right).
\end{equation}

Аэродинамическое демпфирование по тангажу и рысканию вычисляется методом плоских
сечений:
\begin{equation}
\vect{M}_{\text{д}} = -\tfrac12\rho V C_{dc} D\, \mathcal{J}\,
 \boldsymbol{\omega}_\perp,
\qquad
\mathcal{J}=\int_0^L (z-z_{\text{цм}})^2\,\dd z
 = \frac{L^3}{3}-z_{\text{цм}}L^2+z_{\text{цм}}^2 L,
\end{equation}
где $C_{dc}=0{,}6$ — эффективный коэффициент поперечного обтекания.
Демпфирование по крену ввиду большого удлинения корпуса не учитывается.

\subsection{Решётчатые рули}

Четыре руля расположены по азимутам $\phi_i=45^\circ+90^\circ i$ в точках
$\vect{r}_i=\left(R\cos\phi_i,\ R\sin\phi_i,\ z_f\right)$. Местная скорость на
руле учитывает вращение корпуса:
\begin{equation}
\vect{v}_i = \vect{v}_{\text{о},B} + \boldsymbol{\omega}\times\vect{r}_i,
\qquad
q_i=\tfrac12\rho v_i^2,
\qquad
\alpha_i = \delta_i
 + \arctan\frac{\vect{v}_i\cdot\hat{\vect{e}}_{r,i}}{-v_{i,z}},
\end{equation}
где $\delta_i$ — угол отклонения руля, $\hat{\vect{e}}_{r,i}$ — радиальный орт.
Нормальная сила руля
\begin{equation}
L_i = q_i S_f\, C_{N\alpha}^f(M)\,
 \operatorname{clip}\left(\alpha_i,\ \pm\alpha_{\text{ср}}\right) k_{\text{э}},
\end{equation}
с учётом снижения эффективности в трансзвуковой области, где течение в ячейках
руля запирается:
\begin{equation}
C_{N\alpha}^f(M)=C_{N\alpha}^f\cdot
\begin{cases}
1, & M<0{,}6,\\
1-2{,}25\left(M-0{,}6\right), & 0{,}6\le M\le0{,}8,\\
0{,}55, & 0{,}8<M<1{,}2,\\
0{,}55+0{,}3125\left(M-1{,}2\right), & 1{,}2\le M\le 2,\\
0{,}80, & M>2 .
\end{cases}
\end{equation}

% =====================================================================
\section{Алгоритмы наведения атмосферного участка}

\subsection{Тормозной импульс входа}
\label{subsec:entry}

Импульс включается и выключается по прогнозируемому максимальному скоростному
напору. Прогноз выполняется численным интегрированием вертикального снижения с
шагом $\Delta t=1$~с:
\begin{equation}
q_{\max} = \max_i \tfrac12\rho(h_i)v_i^2,
\end{equation}
\begin{equation}
v_{i+1}=v_i+\left[g_0\left(\frac{R_\oplus}{R_\oplus+h_i}\right)^2
 - \frac{\rho(h_i) v_i^2 S C_A^{\text{пр}}}{2m}\right]\Delta t,
\qquad
h_{i+1}=h_i-v_{i+1}\Delta t .
\end{equation}
Импульс включается при $q_{\max}\ge72$~кПа и запасе топлива выше 7~т,
выключается при $q_{\max}\le68$~кПа либо при достижении указанного запаса.

Боковое управление в течение импульса выполняется по методу нулевого промаха по
положению и скорости:
\begin{equation}
a_{\text{б}} = \operatorname{clip}\!\left(
 k_r\frac{-6 r}{t_{\text{ост}}^2} + k_v\frac{-4 v}{t_{\text{ост}}},
 \ \pm a_{\text{дв}}\sin\phi_{\text{кр}}\right),
\qquad
a_{\text{дв}} = \frac{n_e T_e(p)}{m},
\end{equation}
где $r$ и $v$ — боковые отклонение и скорость, $t_{\text{ост}}$ — оценка
оставшегося времени, $\phi_{\text{кр}}=15^\circ$ — предельный угол крена.

\subsection{Посадочный тормозной импульс}

Опорная зависимость вертикальной скорости от высоты соответствует торможению с
постоянным замедлением и выходом на заданную скорость касания $v_{\text{кас}}$:
\begin{equation}
v_{\text{оп}}(h) = -\sqrt{v_{\text{кас}}^2
 + 2 a_{\text{пр}} h_{\text{ост}}},
\qquad
a_{\text{пр}} = k\left(\frac{n_e T_{\text{полн}}}{m}-g_0\right),
\end{equation}
где $k=0{,}85$ при выпущенных рулях и $k=0{,}58$ в противном случае,
$h_{\text{ост}}$ — расстояние от опор до поверхности. Величина $a_{\text{пр}}$
фиксируется в момент включения: её пересчёт по текущей массе приводил бы к
непрерывному увеличению крутизны опорной зависимости по мере выработки топлива и
к отставанию ступени от опорного профиля.

Потребное ускорение, тяга и дросселирование:
\begin{equation}
a_{\text{к}} = a_{\text{пр}} + K_v\left(v_{\text{оп}}-v_z\right),
\qquad
T_{\text{потр}} = m\left(g_0+a_{\text{к}}\right) - D,
\qquad
\tau = \operatorname{clip}\!\left(\frac{T_{\text{потр}}}{n_e T_e(p)},\
 \tau_{\min},\ 1\right),
\end{equation}
где $D$ — сила аэродинамического сопротивления. Боковое управление задаёт
скорость, направленную к точке посадки, пропорционально отклонению, с
ограничением по профилю торможения
$v_{\text{ж}}=\sqrt{2a_{\text{зам}}r}$, что обеспечивает апериодическое
приближение без перерегулирования.

\subsection{Контур управления ориентацией}

Требуемое направление продольной оси определяется вектором потребного ускорения:
\begin{equation}
\hat{\vect{z}}_{\text{ж}} = \frac{\left(a_{\text{б},x},\ a_{\text{б},y},\
 a_{\text{в}}\right)}{\left\lVert\cdot\right\rVert},
\qquad
\vect{e} = \Rot^{-1}(\quat{q})
 \left[\Rot(\quat{q})\hat{\vect{z}}_B\times\hat{\vect{z}}_{\text{ж}}\right],
\end{equation}
причём боковое ускорение предварительно ограничивается конусом допустимого
отклонения продольной оси:
\begin{equation}
a_{\text{б}}^{\max}=a_{\text{в}}\tan\theta_{\max},
\qquad
\theta_{\max}=\operatorname{clip}\!\left(4^\circ
 +\left(\theta_{\text{в}}-4^\circ\right)\frac{h-20}{180},\
 4^\circ,\ \theta_{\text{в}}\right),
\end{equation}
где $\theta_{\text{в}}$ равно $15^\circ$ при выпущенных рулях и $12^\circ$ в
противном случае. Уменьшение допустимого отклонения с уменьшением высоты
обеспечивает касание поверхности всеми опорами одновременно.

Потребный управляющий момент реализуется отклонением камер при $T>5$~кН,
решётчатыми рулями при $q>200$~Па и выпущенных рулях, в остальных случаях — РСУ.
Углы отклонения камер:
\begin{equation}
\sin\delta_0 = \operatorname{clip}\!\left(
 \frac{-M_{y}}{z_{\text{цм}}T},\ \pm\sin\delta_{\max}\right),
\qquad
\sin\delta_1 = \operatorname{clip}\!\left(
 \frac{M_{x}}{z_{\text{цм}}T},\ \pm\sin\delta_{\max}\right).
\end{equation}

% =====================================================================
\section{Контактное взаимодействие с посадочной поверхностью}

Для каждой из четырёх опор при внедрении $\delta_i$, ходе смятия
энергопоглотителя $c_i$ и вертикальной скорости точки контакта $v_z$ упругая и
нормальная составляющие реакции равны
\begin{equation}
F_{\text{у}} = K\left(\delta_i - c_i\right),
\qquad
F_n = \begin{cases}
F_{\text{у}} - C v_z, & F_{\text{у}}\le F_{\text{см}},\\[1ex]
F_{\text{см}} - C v_z,\quad \dot c_i = -v_z,
 & F_{\text{у}}>F_{\text{см}},\ c_i<s_{\max},\\[1ex]
F_{\text{см}} + 10K\left(\delta_i-c_i\right) - C v_z, & c_i\ge s_{\max},
\end{cases}
\end{equation}
где $K$ — жёсткость, $C$ — коэффициент демпфирования, $F_{\text{см}}$ — усилие
смятия энергопоглотителя, $s_{\max}$ — предельный ход смятия. Три строки
соответствуют упругому деформированию, режиму смятия при постоянном усилии и
работе опоры после исчерпания хода энергопоглотителя.

Касательная составляющая определяется законом сухого трения:
\begin{equation}
\vect{F}_t = -\mu F_n \hat{\vect{v}}_t .
\end{equation}

Скорость касания определяется как полная скорость нижней точки опоры в момент
первого контакта и фиксируется с частотой интегрирования. Посадка считается
успешной при скорости касания не выше 6~м/с. Разрушение конструкции фиксируется
при превышении скоростным напором значения 80~кПа в течение более 2~с.

% =====================================================================
\section{Наведение посадочного участка методом выпуклой оптимизации}
\label{sec:pdg}

Наряду с алгоритмом посадочного торможения реализован алгоритм наведения,
основанный на решении задачи оптимального управления методом последовательной
выпуклой аппроксимации. Атмосфера в этой постановке не учитывается: движение
определяется тягой и тяготением.

Область применимости алгоритма ограничена посадочным торможением, то есть
последними километрами траектории. К этому моменту ступень практически пуста:
из стартовой массы $208{,}9$~т остаётся около $32{,}6$~т, из которых $25{,}6$~т
приходится на сухую конструкцию (табл.~\ref{tab:mass}). Поэтому массы в
табл.~\ref{tab:par2} составляют десятки, а не сотни тонн — это не противоречие с
табл.~\ref{tab:par}, а иной участок полёта.

Без атмосферы всё торможение обеспечивается только двигательной установкой,
вследствие чего располагаемая характеристическая скорость ограничивает высоту
начала участка. При параметрах табл.~\ref{tab:par2}
$\Delta v_{\text{расп}}=448$~м/с, что соответствует началу торможения на высоте
порядка 2\dots3~км. Участки входа и аэродинамического спуска, где основное
торможение создаётся сопротивлением воздуха, этим алгоритмом не описываются.

\subsection{Уравнения движения}

Вектор состояния имеет размерность 14, вектор управления — 3:
\begin{equation}
\vect{x}=\left(m,\ \vect{r}_I,\ \vect{v}_I,\ \quat{q},\
 \boldsymbol{\omega}_B\right),
\qquad
\vect{u}=\vect{T}_B .
\end{equation}
\begin{align}
\dot m &= -\frac{\lVert\vect{u}\rVert}{I_{sp}\,g_0}, \label{eq:pdgm}\\
\dot{\vect{r}} &= \vect{v}, \\
\dot{\vect{v}} &= \frac{\Rot(\quat{q})\,\vect{u}}{m} + \vect{g}, \\
\dot{\quat{q}} &= \tfrac{1}{2}\,\quat{q}\otimes(0,\boldsymbol{\omega}), \\
\vect{J}\dot{\boldsymbol{\omega}} &= \vect{r}_T\times\vect{u}
  - \boldsymbol{\omega}\times\vect{J}\boldsymbol{\omega}. \label{eq:pdgw}
\end{align}
Тензор инерции $\vect{J}=\operatorname{diag}(J_x,J_y,J_z)$ в этой постановке
принят постоянным. Точка приложения тяги $\vect{r}_T=(0,0,-\ell)$ лежит на
продольной оси, вследствие чего управляющий момент по крену не создаётся и крен
не управляется.

\subsection{Приведение к задаче со свободным конечным временем}

Свободное конечное время вводится преобразованием независимой переменной: задача
формулируется на нормированном интервале $\varsigma\in[0,1]$,
\begin{equation}
\frac{\dd\vect{x}}{\dd\varsigma} = \sigma\,\vect{f}(\vect{x},\vect{u}),
\qquad \sigma \equiv t_f ,
\end{equation}
благодаря чему время полёта становится скалярной переменной задачи.

На каждой итерации выполняется линеаризация в окрестности текущего опорного
решения $\left(\bar{\vect{x}},\bar{\vect{u}},\bar\sigma\right)$:
\begin{equation}
\frac{\dd\vect{x}}{\dd\varsigma} \approx
  \vect{A}\vect{x} + \vect{B}\vect{u} + \vect{S}\sigma + \vect{z},
\end{equation}
\begin{equation}
\vect{A}=\bar\sigma\frac{\partial \vect{f}}{\partial \vect{x}},\quad
\vect{B}=\bar\sigma\frac{\partial \vect{f}}{\partial \vect{u}},\quad
\vect{S}=\vect{f}\left(\bar{\vect{x}},\bar{\vect{u}}\right),\quad
\vect{z}=-\left(\vect{A}\bar{\vect{x}}+\vect{B}\bar{\vect{u}}\right).
\end{equation}
Дискретизация выполняется на $K$ узлах при кусочно-линейной аппроксимации
управления; переходные матрицы вычисляются методом Рунге—Кутты:
\begin{equation}
\vect{x}_{k+1} = \bar{\vect{A}}_k\vect{x}_k
 + \bar{\vect{B}}^-_k\vect{u}_k + \bar{\vect{B}}^+_k\vect{u}_{k+1}
 + \bar{\vect{S}}_k\sigma + \bar{\vect{z}}_k + \boldsymbol{\nu}_k ,
\end{equation}
где $\boldsymbol{\nu}_k$ — виртуальное управление, обеспечивающее совместность
линеаризованной подзадачи на каждой итерации.

\subsection{Целевая функция}

\begin{equation}
\min\;\; -\frac{m_N}{m_{\text{нач}}}
 \;+\; w_\nu\left(\sum_{k}\sum_{i}\frac{\lvert\nu_{k,i}\rvert}{s_{x_i}}
 \;+\; \sum_{k}\sum_{s}\zeta_{k,s}\right),
\end{equation}
что соответствует максимуму конечной массы, то есть минимуму расхода топлива,
при штрафе на невязку уравнений динамики и на ослабляющие переменные $\zeta$.

\subsection{Ограничения}

Ограничения накладываются в каждом узле и представимы конусом второго порядка:
\begin{align}
m_k &\ge m_{\text{сух}} && \text{(запас топлива)}, \\
\lVert\vect{u}_k\rVert &\le T_{\max} && \text{(верхнее ограничение тяги)}, \\
\bar{\vect{n}}_k^{\mathsf{T}}\vect{u}_k &\ge T_{\min},
  \quad \bar{\vect{n}}_k=\frac{\bar{\vect{u}}_k}{\lVert\bar{\vect{u}}_k\rVert}
  && \text{(нижнее, линеаризованное)}, \\
\lVert\vect{u}_k\rVert &\le \frac{u_{k,z}}{\cos\delta_{\max}}
  && \text{(конус отклонения камер)}, \\
\left\lVert(q_1,q_2)_k\right\rVert
 &\le \sqrt{\tfrac{1}{2}\left(1-\cos\theta_{\max}\right)}
  && \text{(отклонение оси, из \eqref{eq:tilt})}, \\
\lVert\boldsymbol{\omega}_k\rVert &\le \omega_{\max}
  && \text{(угловая скорость)}, \\
\left\lVert \vect{r}_{xy,k}-\vect{r}_{f,xy}\right\rVert
  &\le \operatorname{ctg}\gamma_{\text{г}}
   \left(r_{z,k}-r_{f,z}\right)
  && \text{(конус подхода)} .
\end{align}
Нижнее ограничение тяги является единственным существенно невыпуклым; оно
линеаризуется относительно опорного направления.

Ограничение, включаемое по состоянию, задаётся парой функций $(g,c)$ и выражает
импликацию $g<0\Rightarrow c\le0$ одним непрерывным неравенством
\begin{equation}
-\min(g,0)\cdot c \le 0 .
\end{equation}
Реализованный случай ограничивает угол атаки выше пороговой скорости; при
$\hat{\vect{z}}_B=\Rot(\quat{q})\vect{e}_3$
\begin{equation}
g = v_{\text{пор}} - \lVert\vect{v}\rVert,
\qquad
c = \vect{v}\cdot\hat{\vect{z}}_B
 + \cos\alpha_{\max}\lVert\vect{v}\rVert .
\end{equation}

\subsection{Доверительная область}

Доверительная область задаётся конусным ограничением на масштабированное
отклонение в каждом узле:
\begin{equation}
\left\lVert\left(\hat{\vect{x}}_k-\hat{\bar{\vect{x}}}_k;\
 \hat{\vect{u}}_k-\hat{\bar{\vect{u}}}_k\right)\right\rVert \le r,
\qquad
\left\lvert \hat\sigma-\hat{\bar\sigma}\right\rvert \le r .
\end{equation}
Шаг принимается по отношению действительного снижения штрафной функции к
прогнозируемому:
\begin{equation}
\varrho = \frac{J\left(\bar{\vect{x}}\right)
 - J_{\text{н}}\left(\vect{x}^\star\right)}
 {J\left(\bar{\vect{x}}\right)
 - J_{\text{л}}\left(\vect{x}^\star\right)} .
\end{equation}
Шаг принимается при $\varrho\ge\varrho_0$; радиус уменьшается в $c_-$ раз при
$\varrho<\varrho_1$ и увеличивается в $c_+$ раз при $\varrho\ge\varrho_2$ с
ограничением диапазоном $\left[r_{\min},r_{\max}\right]$. Нелинейная штрафная
функция $J_{\text{н}}$ вычисляется интегрированием принятого управления через
исходные уравнения движения, вследствие чего вычисляемая невязка является
действительной погрешностью интегрирования.

\subsection{Оценка достижимости}

Перед решением выполняется энергетическая оценка:
\begin{equation}
\Delta v_{\text{расп}} = I_{sp} g_0
 \ln\frac{m_{\text{нач}}}{m_{\text{сух}}},
\qquad
v_{\text{пр}} = \sqrt{\lVert\vect{v}_0\rVert^2 + 2 g_0 h_0},
\qquad
\frac{T}{W} = \frac{T_{\max}}{m_{\text{нач}} g_0} .
\end{equation}
Величина $v_{\text{пр}}$ есть скорость, которую набрал бы аппарат при
неработающих двигателях, и является строгой нижней оценкой потребной
характеристической скорости. Задача признаётся неразрешимой при
$v_{\text{пр}}>\Delta v_{\text{расп}}$.

\subsection{Возмущения и контур слежения}

Решение задачи наведения вычисляется однократно, поскольку его получение занимает
время порядка половины секунды и несовместимо с работой в контуре управления;
отслеживание опорной траектории выполняет пропорционально-дифференциальный
регулятор.

Возмущения вводятся в уравнения движения и не сообщаются алгоритмам наведения и
управления: масштаб тяги $s_T$, масштаб удельного импульса $s_I$, перекос камеры
на постоянный угол $a$ вокруг оси с азимутом $\psi$ и запаздывание привода.
Перекос описывается формулой Родрига при
$\hat{\vect{k}}=(-\sin\psi,\cos\psi,0)$:
\begin{equation}
\vect{R} = \vect{I}\cos a + \left[\hat{\vect{k}}\right]_\times \sin a
         + \hat{\vect{k}}\hat{\vect{k}}^{\mathsf{T}}(1-\cos a),
\end{equation}
запаздывание — апериодическим звеном с постоянной времени $\tau_a$:
\begin{equation}
\vect{u} \leftarrow \vect{u} + \left(1-e^{-\Delta t/\tau_a}\right)
  \left(\vect{u}_{\text{ц}}-\vect{u}\right).
\end{equation}

Регулятор формирует ограниченную поправку по ускорению и потребную силу:
\begin{equation}
\vect{a}_c = \operatorname{sat}_{a_{\max}}\!\left[
  k_{p}\left(\vect{r}_r-\vect{r}\right)
  + k_{d}\left(\vect{v}_r-\vect{v}\right)\right],
\qquad
\vect{f}_{\text{ж}} = \Rot\left(\quat{q}_r\right)\vect{u}_{r,B}
 + m\,\vect{a}_c ,
\end{equation}
при $\lVert\vect{f}_{\text{ж}}\rVert\in\left[T_{\min},T_{\max}\right]$. Требуемая
ориентация получается композицией опорной ориентации с малым поворотом,
переводящим $\hat{\vect{f}}_r$ в $\hat{\vect{f}}_{\text{ж}}$:
\begin{equation}
\vect{a}_{\text{п}}=\hat{\vect{f}}_r\times\hat{\vect{f}}_{\text{ж}},
\qquad
\dd\quat{q}=\operatorname{norm}
 \left(1,\ \tfrac{1}{2}\vect{a}_{\text{п}}\right),
\qquad
\quat{q}_{\text{ж}}=\dd\quat{q}\otimes\quat{q}_r .
\end{equation}
Ошибка ориентации
$\quat{q}_e=\quat{q}^{-1}\otimes\quat{q}_{\text{ж}}$ приводится к двум
управляемым составляющим, после чего формируются угловые ускорения и приращения
управления:
\begin{gather}
e_x = 2\operatorname{sign}\left(q_{e,0}\right)q_{e,1},
\qquad
e_y = 2\operatorname{sign}\left(q_{e,0}\right)q_{e,2},
\qquad
a_x = k_{a}e_x + k_{d,a}\left(\omega_{r,x}-\omega_x\right),\\[4pt]
u_x = u_{r,x}\frac{\lVert\vect{f}_{\text{ж}}\rVert}{\lVert\vect{f}_r\rVert}
 - \frac{J_y a_y}{L},
\qquad
u_y = u_{r,y}\frac{\lVert\vect{f}_{\text{ж}}\rVert}{\lVert\vect{f}_r\rVert}
 + \frac{J_x a_x}{L},\\[4pt]
u_z=\sqrt{\max\!\left(\lVert\vect{f}_{\text{ж}}\rVert^2-u_x^2-u_y^2,\;
  \tfrac{1}{4}\lVert\vect{f}_{\text{ж}}\rVert^2\right)},
\end{gather}
где $L=\lvert r_{T,z}\rvert$ — плечо приложения тяги. Поперечная составляющая
ограничена половиной командного модуля.

% =====================================================================
\section{Восстановление параметров по сокращённой телеметрии}

Расчёт атмосферного участка регистрируется с частотой 2~Гц в виде набора
\begin{equation}
\left\{t,\ h,\ v_z,\ \tau,\ \lvert\theta\rvert,\
 \ell,\ v_{\text{рад}},\ q,\ \omega_\perp,\ w,\ m,\ \text{режим}\right\},
\end{equation}
где $\ell=\sqrt{r_x^2+r_y^2}$ — расстояние до точки посадки. Для визуализации
требуются три величины, отсутствующие в этом наборе.

\subsection{Тяга по секундному расходу}

Степень дросселирования не определяет числа работающих камер, поэтому тяга
вычисляется по расходу:
\begin{equation}
T = \max\!\left(-\frac{\dd m}{\dd t},\ 0\right) g_0\, I_{sp}(h),
\qquad
I_{sp}(h) = I_{\text{п}}
 + \left(I_{\text{з}}-I_{\text{п}}\right) e^{-h/H},
\end{equation}
при $H=8500$~м; тяга полагается нулевой при нулевом дросселировании.

\subsection{Знак дальности}

Величина $\ell$ неотрицательна, поэтому траектория, проходящая над точкой
посадки, отображается зеркально отражённой. Признаком прохождения служит
одновременное выполнение трёх условий:
\begin{equation}
v_{\text{рад},i-1}\,v_{\text{рад},i} < 0,
\qquad
\min\left(\left\lvert v_{\text{рад},i-1}\right\rvert,
 \left\lvert v_{\text{рад},i}\right\rvert\right) > v_\epsilon,
\qquad
\min\left(\ell_{i-1},\ell_i\right)
 < \max_i\left\lvert v_{\text{рад}}\right\rvert\Delta t ,
\end{equation}
при $v_\epsilon=0{,}5$~м/с, после чего знак инвертируется для всех последующих
отсчётов. Действительный разворот по радиусу характеризуется плавным переходом
$v_{\text{рад}}$ через нуль и не удовлетворяет двум последним условиям. Операция
выполняется до передискретизации, поскольку знакопеременная величина проходит
через нуль гладко и допускает линейную интерполяцию.

\subsection{Направление отклонения продольной оси}

В телеметрии фиксируется модуль отклонения, но не его направление. Тяга
действует вдоль продольной оси, поэтому направление определяется знаком бокового
ускорения:
\begin{equation}
\operatorname{sign}(\theta)
 = \operatorname{sign}\left(\overline{\dot v_{\text{рад}}}\right)
\quad\text{при } T>0{,}05\,T_{\max},
\qquad
\theta_{\text{к}} = \lvert\theta\rvert\operatorname{sign}(\theta),
\end{equation}
причём ускорение сглаживается пятиточечным скользящим средним, а при выключенных
двигателях сохраняется последнее определённое значение.

\subsection{Ограничение угловой скорости}

Полученная зависимость $\theta_{\text{к}}(t)$ разрывна в точках смены знака.
Располагаемое угловое ускорение определяется суммой трёх управляющих моментов:
\begin{equation}
M_{\text{к}} = T \sin\delta_{\max}\; z_{\text{цм}},
\quad
M_{\text{р}} = n_f\, q\, S_f\, C_{N\alpha}^f\, \delta_f
  \left(z_f - z_{\text{цм}}\right),
\quad
M_{\text{РСУ}} = n_r F_r \left(z_r - z_{\text{цм}}\right),
\end{equation}
\begin{equation}
\varepsilon_{\max}
 = \frac{M_{\text{к}}+M_{\text{р}}+M_{\text{РСУ}}}{J},
\qquad
t_{\text{р}} = 2\sqrt{\frac{\theta}{\varepsilon_{\max}}},
\qquad
\omega_{\max} = \sqrt{\theta\,\varepsilon_{\max}} .
\label{eq:eps}
\end{equation}
Каждый источник преобладает на своём участке полёта: рули — на аэродинамическом
спуске, отклонение камер — на участках работы двигателей, РСУ — в разрежённой
атмосфере.

Слежение выполняется оптимальным по быстродействию законом:
\begin{equation}
\omega_{\text{ж}} = \operatorname{sign}(e)
 \sqrt{2\varepsilon_{\max}\lvert e\rvert},
\qquad e = \theta_{\text{к}} - \theta,
\end{equation}
\begin{equation}
\omega \leftarrow \operatorname{clip}\Bigl(
  \omega + \operatorname{clip}\left(\omega_{\text{ж}}-\omega,\
   \pm\varepsilon_{\max}\Delta t\right),
  \ \pm\omega_\perp\Bigr),
\qquad
\theta \leftarrow \theta + \omega\,\Delta t .
\end{equation}
Ограничение зарегистрированным значением $\omega_\perp$ является точным:
составляющая угловой скорости в плоскости не может превышать модуля этой
скорости.

\subsection{Восстановленная ориентация}

\begin{equation}
\quat{q} = \left(\cos\frac{\theta}{2},\ 0,\
 \sin\frac{\theta}{2},\ 0\right),
\qquad
\boldsymbol{\omega} = \left(0,\ \dot\theta,\ 0\right),
\qquad
\vect{u}=(0,0,T).
\end{equation}
Запись с частотой 2~Гц интерполируется линейно на равномерную сетку 20~Гц для
воспроизведения; метка режима полёта не интерполируется, а сохраняется до
следующего отсчёта.

% =====================================================================
\section{Проекция для визуализации}

Ось дальности совпадает с горизонтальным направлением начального положения, ось
бокового отклонения дополняет систему:
\begin{equation}
\eh = \frac{\vect{r}_{0,xy}}{\lVert\vect{r}_{0,xy}\rVert},
\qquad
\ec = \hat{\vect{z}}\times\eh,
\qquad
d=\vect{r}\cdot\eh,\quad c=\vect{r}\cdot\ec,
\qquad
\hat{\vect{b}}_I = \Rot(\quat{q})\vect{e}_3 .
\end{equation}
Производные углы:
\begin{align}
\theta &= \arccos\left(\operatorname{clip}
 \left(1-2\left(q_1^2+q_2^2\right),-1,1\right)\right), \\
\delta &= \operatorname{arctg}
 \frac{\lVert(u_x,u_y)\rVert}{\max(u_z,\epsilon)}, \\
\alpha &= \arccos\!\left(\operatorname{clip}\!\left(
  \frac{-\vect{v}}{\lVert\vect{v}\rVert}\cdot\hat{\vect{b}}_I,-1,1\right)\right).
\end{align}

Масштаб изображения постоянен и одинаков по обеим осям. Переход от обзора всей
траектории к сопровождению аппарата выполняется геометрической интерполяцией
масштаба и сглаженной по времени интерполяцией положения:
\begin{equation}
\tilde s = s^2(3-2s),
\qquad
w(s) = W\left(\frac{w_c}{W}\right)^{\tilde s},
\qquad
\vect{p}(s) = (1-\tilde s)\,\vect{p}_{\text{сц}}
 + \tilde s\,\vect{p}_{\text{ап}} .
\end{equation}
Переход инициируется по бегущему минимуму высоты, отсчитываемому от последнего
отсчёта выше порогового значения, а не от начала полёта; в противном случае для
траектории, начинающейся на поверхности, условие оказывается выполненным до
старта. По тому же правилу определяется момент выпуска посадочных опор.

% =====================================================================
\section{Параметры модели}

\begin{table}[htb]
\centering
\caption{Параметры ступени}
\label{tab:par}
\begin{tabular}{llr}
\toprule
Обозначение & Наименование & Значение \\
\midrule
$L$ & Полная длина & $47{,}7$ м \\
$L_{\text{ст}}$ & Длина без межступенного отсека & $41{,}2$ м \\
$D$; $R$ & Диаметр; радиус & $3{,}66$; $1{,}83$ м \\
$S$ & Площадь миделя & $10{,}52\ \text{м}^2$ \\
$m_{\text{сух}}$ & Сухая масса с опорами и рулями & $25\,600$ кг \\
$z_{\text{сух}}$ & Центр масс сухой конструкции & $12{,}4$ м \\
$A$; $z_{\text{осн}}$; $\rho_{\text{т}}$ & Параметры топливного столба
 & $9{,}9\ \text{м}^2$; $1{,}6$ м; $1048\ \text{кг/м}^3$ \\
$T_{\text{з}}$; $T_{\text{п}}$ & Тяга камеры (земная; пустотная)
 & $845$; $931{,}9$ кН \\
$I_{\text{з}}$; $I_{\text{п}}$ & Удельный импульс & $282$; $311$ с \\
$\dot m_e$; $A_e$ & Секундный расход; эффективная площадь среза
 & $305{,}6$ кг/с; $0{,}859\ \text{м}^2$ \\
$\tau_{\min}$ & Нижний предел дросселирования & $0{,}40$ \\
$\delta_{\max}$ & Предельный угол отклонения камер & $5^\circ$ \\
$n_f$; $S_f$; $C_{N\alpha}^f$ & Число рулей; площадь; производная
 & $4$; $2{,}4\ \text{м}^2$; $3{,}0$ 1/рад \\
$\delta_f$; $z_f$ & Предельное отклонение; координата рулей
 & $20^\circ$; $45{,}0$ м \\
$n_r$; $F_r$; $z_r$ & Число сопел; тяга; координата РСУ
 & $2$; $500$ Н; $40{,}5$ м \\
$K$; $C$ & Жёсткость и демпфирование опоры
 & $3{\cdot}10^6$ Н/м; $1{,}2{\cdot}10^5$ Н$\cdot$с/м \\
$F_{\text{см}}$; $s_{\max}$ & Усилие и ход смятия
 & $4{\cdot}10^5$ Н; $0{,}40$ м \\
$\mu$ & Коэффициент трения & $0{,}45$ \\
$q_{\max}$ & Предельный скоростной напор & $80$ кПа \\
\bottomrule
\end{tabular}
\end{table}

\begin{table}[htb]
\centering
\caption{Характерные массы по участкам полёта}
\label{tab:mass}
\begin{tabular}{lrr}
\toprule
Момент & Масса, т & Топливо, т \\
\midrule
Старт & $208{,}9$ & $183{,}3$ \\
Выключение маршевых двигателей & $86{,}59$ & $60{,}99$ \\
Граница входа ($h_{\text{гр}}=65$ км) & $55{,}39$ & $29{,}79$ \\
Начало посадочного торможения ($h=2{,}7$ км) & $32{,}60$ & $7{,}00$ \\
Касание & $27{,}62$ & $2{,}02$ \\
\bottomrule
\end{tabular}

\smallskip
\footnotesize Примечание: сухая масса ступени $25{,}6$~т. Для сравнения,
полностью заправленная ступень имеет стартовую массу $433{,}1$~т при запасе
топлива $410{,}9$~т; рассматриваемый автономный полёт не требует полной
заправки, и потребный запас топлива является результатом решения задачи
выведения, а не исходной величиной.
\end{table}

\begin{table}[htb]
\centering
\caption{Параметры задачи наведения посадочного участка}
\label{tab:par2}
\begin{tabular}{llr}
\toprule
Обозначение & Наименование & Значение \\
\midrule
$m_{\text{нач}}$; $m_{\text{сух}}$ & Масса в начале торможения; сухая масса
 & $30\,100$; $25\,600$ кг \\
$I_{sp}$ & Удельный импульс (земной) & $282$ с \\
$T_{\min}$; $T_{\max}$ & Пределы тяги (одна камера) & $338$; $845$ кН \\
$T_{\min}/G$ & Тяговооружённость при сухой массе & $1{,}35$ \\
$\vect{J}$ & Тензор инерции
 & $(5{,}30;\,5{,}30;\,0{,}093)\cdot10^6\ \text{кг}\cdot\text{м}^2$ \\
$\vect{r}_T$ & Точка приложения тяги & $(0;\,0;\,-9{,}4)$ м \\
$\theta_{\max}$; $\delta_{\max}$; $\gamma_{\text{г}}$ & Предельные углы
 & $30^\circ$; $5^\circ$; $20^\circ$ \\
$\omega_{\max}$ & Предельная угловая скорость & $20^\circ$/с \\
$K$; $w_\nu$ & Число узлов; вес виртуального управления & $31$; $10^3$ \\
\bottomrule
\end{tabular}

\smallskip
\footnotesize Примечание: запас топлива на посадочное торможение ($4{,}5$~т)
получен из формулы Циолковского для потребной характеристической скорости около
400~м/с при сухой массе $25{,}6$~т. Поскольку $T_{\min}/G>1$, зависание
невозможно, и торможение выполняется единственным включением без возможности
исправления ошибки момента запуска.
\end{table}

% =====================================================================
\section{Принятые допущения}

1. Вращение Земли и боковое движение не учитываются. При пуске в восточном
направлении с широты $28{,}5^\circ$ начальная переносная скорость составляет
408~м/с и возникает кориолисово боковое отклонение. Модель внеатмосферного
участка плоская, что вносит погрешность в характеристическую скорость манёвра
возврата и в положение трассы; возникающая невязка устраняется коррекцией точки
прицеливания, а не учитывается непосредственно.

2. На участке выведения угловое движение не моделируется: задаётся направление
тяги, которое считается реализуемым. Скорость разворота на манёвре возврата
ограничена располагаемым управляющим моментом отклонения камер, а на
баллистическом участке — моментом РСУ, однако замкнутый контур управления
ориентацией с моделью привода в этой постановке отсутствует.

3. В задаче наведения посадочного участка (\S\ref{sec:pdg}) тензор инерции принят
постоянным, тогда как за время торможения расходуется около $15\,\%$ массы
аппарата. На атмосферном участке такого допущения не делается: там тензор
инерции пересчитывается по \eqref{eq:inertia}.

4. Поперечное ограничение в контуре слежения (\S\ref{sec:pdg}) соответствует
отклонению камер на $30^\circ$, что существенно превышает конструктивный предел
$\delta_{\max}=5^\circ$, принятый при решении задачи наведения. Ограничение
действует только в контуре управления, однако реализуемая траектория не
гарантирует соблюдения конуса отклонения.

5. Отсутствие возможности зависания ($T_{\min}/G=1{,}35$ при сухой массе) делает
посадочное торможение чувствительным к моменту запуска двигателя: однократно
вычисленная опорная траектория и следящий регулятор не обеспечивают компенсации
отклонения тяги порядка $2\,\%$. Устранение этого недостатка требует перехода к
повторному решению задачи наведения на скользящем горизонте.

6. Нижнее ограничение тяги линеаризуется относительно опорного решения,
вследствие чего допустимое множество зависит от текущей итерации.

7. Управляющий момент решётчатых рулей в \eqref{eq:eps} оценивается сверху:
принимается предельное отклонение при текущем скоростном напоре без учёта срыва
потока и затенения.

8. Обработка телеметрии предполагает, что горизонтальное движение происходит в
одной вертикальной плоскости; азимут в регистрируемом наборе отсутствует, и
проверка этого предположения по данным невозможна.

9. Устойчивость посадки существенно зависит от ветрового воздействия. При средней
скорости ветра 8~м/с и высокой интенсивности турбулентности одна и та же
траектория приводит либо к скорости касания $2{,}8$~м/с, либо к разрушению
конструкции при скорости $27{,}6$~м/с в зависимости только от заправки. Основной
расчётный случай принят при скорости ветра 3~м/с.

10. Крен не управляется и при проекции для визуализации не учитывается.

\end{document}
